%%%
%%% Radical "⽑"
%%%
\section*{Radical 82: ``⽑''}\addcontentsline{toc}{section}{Radical 82: ⽑}\addcontentsline{loh}{figure}{\#\#\#\# 82: ⽑}

%%%%%%%%%% 毛 %%%%%%%%%%
\subsection*{毛}\addcontentsline{loh}{figure}{毛}

\begin{Entry}{毛}{4}{⽑}[Kangxi 82]
  \begin{Phonetics}{毛}{mao2}[][HSK 1,3]
    \definition*{s.}{Sobrenome: Mao}
    \definition{adj.}{bruto; semiacabado | grosseiro | pequeno | fino | descuidado; rude; precipitado | assustado; nervoso; em pânico | impetuoso | rústico; sem acabamento | impuro | (moeda) que não vale mais seu valor nominal; depreciado}
    \definition{clas.}{mao, uma unidade fracionária de dinheiro na China; dez centavos; uma peça de dez centavos}
    \definition[根,撮,身]{s.}{(animal, planta, etc.) cabelo; pena; penugem | (humanos) cabelo; barba | planta; colheita | lã | mofo; bolor}
    \definition{v.}{depreciar; desvalorizar; refere"-se à desvalorização da moeda | (cavalos, gado, etc.) assustar"-se; sentir medo}
  \synonymref{估}{gu4}
  \synonymref{慌}{huang1}
  \synonymref{惊}{jing1}
  \synonymref{怒}{nu4}
  \synonymref{细}{xi4}
  \synonymref{小}{xiao3}
  \antonymref{大}{da4}
  \end{Phonetics}
\end{Entry}

\begin{Entry}{毛巾}{4,3}{⽑,⼱}
  \begin{Phonetics}{毛巾}{mao2jin1}[][HSK 4]
    \definition[条,块]{s.}{toalha; toalha de banho}
  \end{Phonetics}
\end{Entry}

\begin{Entry}{毛衣}{4,6}{⽑,⾐}
  \begin{Phonetics}{毛衣}{mao2yi1}[][HSK 4]
    \definition[件,个]{s.}{suéter; blusa feita de lã}
  \end{Phonetics}
\end{Entry}

\begin{Entry}{毛病}{4,10}{⽑,⽧}
  \begin{Phonetics}{毛病}{mao2bing4}[][HSK 3]
    \definition[个,点,种,些]{s.}{doença ou deficiência; condição de saúde precária ou deficiência física | problema; fracasso; onde o produto está com defeito ou não funciona corretamente | mau hábito; deficiência; falhas no comportamento humano}
  \seealsoref{故障}{gu4zhang4}
  \seealsoref{缺点}{que1dian3}
  \synonymref{弊端}{bi4duan1}
  \synonymref{差错}{cha1cuo4}
  \synonymref{错误}{cuo4wu4}
  \synonymref{坏处}{huai4chu5}
  \synonymref{缺陷}{que1xian4}
  \synonymref{问题}{wen4ti2}
  \synonymref{障碍}{zhang4'ai4}
  \antonymref{长处}{chang2chu4}
  \antonymref{健康}{jian4kang1}
  \antonymref{亮点}{liang4dian3}
  \antonymref{完好}{wan2hao3}
  \antonymref{优点}{you1dian3}
  \end{Phonetics}
\end{Entry}

\begin{Entry}{毛笔}{4,10}{⽑,⽵}
  \begin{Phonetics}{毛笔}{mao2bi3}[][HSK 5]
    \definition[支,枝,根,管]{s.}{pincel para escrever; pincel chinês; canetas feitas com pelos de coelho, carneiro, doninha, etc., são materiais tradicionais utilizados para escrever caracteres chineses e pintar pinturas tradicionais chinesas}
  \end{Phonetics}
\end{Entry}

%%%%%%%%%% 毯 %%%%%%%%%%
\subsection*{毯}\addcontentsline{loh}{figure}{毯}

\begin{Entry}{毯}{12}{⽑}
  \begin{Phonetics}{毯}{tan3}
    \definition[条]{s.}{cobertor; tapete; carpete}
  \end{Phonetics}
\end{Entry}

\begin{Entry}{毯子}{12,3}{⽑,⼦}
  \begin{Phonetics}{毯子}{tan3zi5}[][HSK 7-9]
    \definition[条,张,床,面]{s.}{cobertor; tecidos grossos e macios podem ser usados como cortinas, capas ou decorações}
  \synonymref{被子}{bei4zi5}
  \synonymref{垫子}{dian4zi5}
  \end{Phonetics}
\end{Entry}

%%%%% EOF %%%%%

